\RequirePackage{plautopatch}
\documentclass[uplatex,dvipdfmx,a4paper,10pt]{jsarticle}
\usepackage{graphicx}
\usepackage{amsmath}
\usepackage{mathcomp}
\usepackage{csvsimple}
\usepackage{siunitx}
\usepackage{url}


\title{比熱の測定}
\author{電気通信大学 I\hspace{-1.2pt}I\hspace{-1.2pt}I類\\
2313213 山口海音}
\date{\today}

\begin{document}

\maketitle

\section{目的}
先ず，ニュートンの冷却の法則が成り立つことを確認する．次に，ジュール熱の存在を認め，熱力学第一法則がなりたつことを仮定したうえで，水とアルコールの比熱を求める．

\section{原理}

\subsection{ニュートンの冷却の法則}
物体の温度の時間変化率$\dfrac{\Delta T}{\Delta t}$は物体と周囲の温度差$T-T_{\mathrm{env}}$に比例する．これをもって物体ごとに決まる比例定数$k$を定める．これは，粘性抵抗を受けて落下する物体の運動にちょうど等しいので，温度は時間に対して以下のように変化する．
\begin{align}
    \frac{T-T_\mathrm{env}}{T_0-T_\mathrm{env}}=e^{-kt}
\end{align}
この法則を確認するためには，$\ln(T-T_\mathrm{env})$と$t$の関係が$1$次方程式で表せるか，あるいは，近似できるかを確認すればよい．
%ニュートンがdQなど使っていないことを示す文献

\subsection{ジュール熱}
抵抗値$R$の物質に電流$I$を$t$秒間流し続けると，熱$Q=RI^2t$が生じる．

\subsection{熱力学第一法則}
内部エネルギーの変化$\Delta U$は，入ってきた熱量$Q$，内部気体がした仕事$W$の差分に等しい．

\subsection{熱容量}
物質の温度を$\Delta T$だけ上げるのに必要な熱量$Q$は$\Delta T$に比例する．これをもって物体ごとに決まる比例定数$L$を定める．
% 熱量とは
% 何故比例定数をそう入れたか
$C$は体積に比例する，即ちその物質の質量に比例する．これをもって物質ごとに決まる比例定数$c$を定める．
これを実験的に求めるには，ある物体に電流を加えて生じさせたジュール熱を全て物質の加熱に使えばよい．即ち，以下が成り立つ．
\begin{align}
    C\Delta T = Q = RI^2t
\end{align}


\section{実験方法}
\subsection{比熱の測定}
第1に，抵抗線を作成し，抵抗値を測定する．
第2に，凡そ$RI^2 = 12\mathrm{W} $となるように電流$I$を設定し，測定する．
第3に，熱量計と用意した水の質量を測定する．
第4に，恒温槽内にスペーサと水を入れた熱量計と温度計を入れる．
第5に，ゆっくり攪拌し，容器内の温度を一定に保．
第6に，電流を流し始め，温度変化$\Delta T = 0.3℃$毎に時間を記録する．
第7に，温度変化$\Delta T = 15℃$ほどになった段階で，記録を終了する．
第8に，残った水の質量を測定する．

次に，水をエタノールに変えて 1 - 7 について同様の実験をする．

\subsection{ニュートンの冷却の法則}
第1に，前項で加熱したエタノールを，電流は止めてそのまま攪拌を続けて記録する．
第2に，温度変化$\Delta T = -10℃$ほどになった段階で，記録を終了する．
第3に，残ったエタノールの質量を測定する．

\section{実験結果}
抵抗線の抵抗値，設定した電流，熱量計の比熱，熱量計の質量，水の質量，エタノールの質量を表\ref{tb:const}に示す．

\begin{table}[htbp]
    \begin{center}
        \caption{比熱を求めるために必要な定数}%naming びみょすぎｗ
        \label{tb:const}
        \begin{tabular}{r|r}
            \hline
            データ名 & 測定値\\
            \hline
            抵抗値 & $(4.413 \pm 0.004)\mathrm{Ω}$\\
            電流 & $(1.675 \pm 0.007)\mathrm{A}$\\
            熱量計の比熱 & $0.383 \mathrm{Jg^{-1}K^{-1}}$\\ %ソースを示せ．
            熱量計の質量 & $(55.06\pm 0.01) \mathrm{g}$\\
            水の質量 & $(179.35 \pm 0.02) \mathrm{g}$\\
            エタノールの質量 & $(139.94 \pm 0.01)\mathrm{g}$\\
            \hline
        \end{tabular}
    \end{center}
\end{table}

水を加熱したときの時間$t$に対する水温$T$の変化を図$1$に示す．

このグラフの傾き$a$と不確かさ$\Delta a$は図$1$のグラフの傾きを読んで以下のようである．

\begin{align}
    a = \frac{T}{t}\\
    \frac{\Delta a}{a} = \sqrt{}\\
\end{align}

\begin{table}[htbp]
    \begin{center}
        \caption{比熱を求めるために必要な定数}%naming びみょすぎｗ
        \label{tb:const}
        \begin{tabular}{r|r}
            \hline
            データ名 & 測定値\\
            \hline
            抵抗値 & $(4.413 \pm 0.004)\mathrm{Ω}$\\
            \hline
        \end{tabular}
    \end{center}
\end{table}

エタノールを加熱したときの時間$t$に対する水温$T$の変化を図$2$に示す．

得られた比熱を表\ref{tb:heatcap}に示す．
\begin{table}[htbp]
    \begin{center}
        \caption{比熱の測定値と文献値}
        \label{tb:heatcap}
        \begin{tabular}{r|rr}
            \hline
            データ名 & 測定値 & 文献値\\
            \hline
            水 &  & \\
            アルコール & & \\
            \hline
        \end{tabular}
    \end{center}
\end{table}

エタノールを冷却したときの時間$t$に対する水温$T$の変化を図$3$に示す．
エタノールを冷却したときの時間$t$に対する水温の変化の対数$\ln(T-T_\mathrm{env})$を図$4$に示す．
ただし，エタノールを冷却したときに攪拌しなかった．また，途中で一度だけ攪拌した．

\section{考察}
\subsection{冷却実験における試料の体積}
\subsection{試料の体積変化による誤差}
\subsection{考慮しきれていない因子}
空気，スペーサ，抵抗線，はんだ，
\subsection{結果の検討}

% \begin{thebibliography}{99}
%     \bibitem{phy} \verb|伊東敏雄，『な～るほど！の力学』，東京学術図書出版社2003．|
%     \bibitem{jihou} \verb|国土地理院，『国土地理院時報』，131集，国土地理院，2018．|
%     \bibitem{sci} \verb|『日本各地の重力実測値』，理科年表プレミアム - コンテンツ表示，理科年表オフィシャルサイト(国立天文台・丸善出版)，2023-07-13閲覧．|
%     \bibitem{keido} \verb|『東京都 市区町村の役所・役場及び東西南北端点の経度緯度(世界測地系)』，国土地理院，2023-07-14閲覧，https://www.gsi.go.jp/common/000230949.pdf．|
% \end{thebibliography}

\end{document}